\chapter{EasyCrypt Framework}
\label{ch:framework}

\section{Modules and Adversaries}

The proof uses EasyCrypt modules to represent schemes, oracles, adversaries,
and reductions.  A module type describes the operations available to a game or
adversary.  A concrete module instantiates those operations.  This is essential
for cryptographic reductions because an adversary is not a function in
isolation; it is code linked against a particular oracle interface.

Mathematically, an EasyCrypt module should be read as an indexed family of
probabilistic transition kernels.  A procedure
\[
  \mathtt{proc}\ f(x):\tau
\]
denotes a subdistribution over return values and final memories.  A module type
specifies which kernels an environment may call.  This is the formal analogue of
the oracle-access notation in cryptographic papers:
\[
  A^{H,\mathsf{Sign}}(pk).
\]
The superscript is not decoration; it is an interface contract.  The EasyCrypt
module system makes that contract explicit and prevents a proof from silently
using state or procedures that the adversary should not possess.

Functor-style modules are the proof mechanism behind reductions.  A reduction is
code that links an adversary against one game to oracles for another game.  In
mathematical notation this is usually written as \(B^O(A)\).  In EasyCrypt it is
a module parameterized by both adversary and oracle modules.  The type checker
therefore enforces which information crosses the reduction boundary.

\section{Probabilistic Judgments}

The proof uses several EasyCrypt judgment forms.  Hoare judgments prove state
properties of programs.  Probabilistic Hoare judgments prove probability bounds.
Equivalence judgments prove that two games produce related outputs under
related states.  Losslessness lemmas establish that procedures terminate with
probability one.

The corresponding mathematical readings are:
\[
\begin{array}{ll}
\{P\}\;c\;\{Q\}
&
\text{for every initial memory satisfying }P,\text{ every terminating final memory satisfies }Q,\\[0.4em]
\Pr_c[Q\mid P]\le \alpha
&
\Pr[c\text{ terminates in a }Q\text{-state}\mid P]\le \alpha,\\[0.4em]
c_1 \sim c_2 : P \Longrightarrow Q
&
\text{there is a coupling of }c_1,c_2\text{ taking }P\text{-related states to }Q\text{-related states}.
\end{array}
\]
The exact EasyCrypt surface syntax differs from this display, but the semantic
content is this combination of state transformers, probability measures, and
relations.

Game-hopping proofs usually require all three forms.  A relational equivalence
shows that two games are identical until a bad event; a probabilistic Hoare
judgment bounds the bad event; ordinary Hoare judgments maintain invariants such
as query-log consistency and budget monotonicity.  This decomposition is the
machine-checked version of the fundamental lemma of game playing.

\section{Fundamental Lemma of Game Playing}

The fundamental lemma can be stated in total-variation language.  Let \(G_0\)
and \(G_1\) be two games with Boolean outputs, and let \(\mathsf{Bad}\) be an
event.  If the games are coupled so that their views and outputs agree whenever
\(\mathsf{Bad}\) does not occur, then
\[
  \left|\Pr[G_0\Rightarrow 1]-\Pr[G_1\Rightarrow 1]\right|
  \le \Pr[\mathsf{Bad}].
\]
Many HAETAE proof obligations are instances of this lemma.  The difficult part
is not the inequality itself; it is proving that the relevant games are truly
identical on the complement of the bad event.  For ROM programming, this means
the programmed point was fresh.  For sampler lifting, this means the sampler
query was not present in the adversary hash log or in the sampler self-log.  For
budgeted wrappers, this means oracle counters and logs remain synchronized.

\section{Proof Style}

The development uses named lemmas heavily.  This makes the proof auditable:
rather than compressing all reasoning into a single top-level statement, it exposes
intermediate facts such as oracle preservation, bad-event bounds, query-budget
invariants, and reduction-term arithmetic.

A useful way to read the proof scripts is as a typed refinement of a
pen-and-paper proof outline.  Paper proofs often use phrases such as ``by
freshness,'' ``by rejection sampling,'' or ``by a standard hybrid argument.''
In the EasyCrypt development each such phrase must correspond to a theorem with
explicit hypotheses.  The main review question is therefore not only whether the
top-level checked implication compiles, but whether every lemma name and
hypothesis corresponds to the intended mathematical invariant.
